\documentclass[letterpaper,11pt]{moderncv}
\moderncvstyle{banking}
\usepackage{xpatch}
\moderncvcolor{black}
\AtBeginDocument{\definecolor{color2}{rgb}{0,0,0}}
\usepackage[letterpaper, margin=0.6in]{geometry}
\usepackage{xcolor}
\usepackage{graphicx}
\usepackage{fontspec}
\setmainfont{Times New Roman}
\pagenumbering{arabic}
\usepackage[unicode]{hyperref}
\sloppy
\usepackage{enumitem}

\usepackage[english]{babel}
\usepackage[autostyle, english = american]{csquotes}
\MakeOuterQuote{"}

%==================== PERSONAL INFO ====================%
\name{\textcolor{black}{Saeyoung}}{Kim}
\address{Cambridge, MA, USA}
\phone[mobile]{+1 (351) 322 5510}
\email{saeyoung\_kim@uml.edu}
\social[linkedin]{saeyoung-macx-kim}
\social[github]{PrimaryAnomaly}

\begin{document}

\makecvtitle
\vspace{-3em}

%==================== SUMMARY ====================%
\section{Summary}
Hands-on technician and engineer with experience assembling, testing, and troubleshooting electromechanical hardware in aerospace environments. Performed precision mechanical assembly and electrical integration of satellite flight hardware in cleanroom facilities. Skilled in following controlled procedures, maintaining documentation, and collaborating with engineering teams to resolve build issues. Background in soldering, PCB rework, test equipment operation, and quality-driven manufacturing practices.

\section{Technical Skills}
\begin{minipage}[t]{0.48\textwidth}
\cvitem{Assembly}{Mechanical Assembly, Electrical Integration, Cleanroom Operations, Precision Hardware Handling}
\cvitem{Electronics}{Soldering, PCBA Rework, Cable/Harness Assembly, PCB Prototyping}
\end{minipage}%
\hfill
\begin{minipage}[t]{0.48\textwidth}
\cvitem{Instruments}{Oscilloscopes, DMMs, Power Supplies, Function Generators, Calipers}
\cvitem{Tools}{Hand Tools, Torque Wrenches, CAD (Fusion360, SolidWorks), MS Office}
\end{minipage}

%==================== EXPERIENCE ====================%
\section{Experience}

\cventry{2025.07--Present}{\textbf{Postdoctoral Associate} | Process Monitoring}{University of Massachusetts Lowell}{Lowell, MA, USA}{}{
\begin{itemize}[leftmargin=1cm, itemsep=0pt, parsep=0pt]
    \item Developed standardized procedures and technical documentation for manufacturing process monitoring
    \item Identified process deviations through data analysis, supporting root cause investigation and corrective actions
\end{itemize}}

\cventry{2024.02--2025.05}{\textbf{Systems Integration Engineer} | Satellite Assembly, Integration, and Testing}{Hanwha Systems}{Seoul, South Korea}{}{
\begin{itemize}[leftmargin=1cm, itemsep=0pt, parsep=0pt]
    \item Performed hands-on mechanical assembly and electrical integration of military satellite hardware in ISO cleanroom environments
    \item Followed and executed controlled assembly and test procedures for satellite subsystem and system-level builds
    \item Operated test equipment (oscilloscopes, DMMs, power supplies) for hardware verification during integration
    \item Troubleshot electromechanical issues during assembly, collaborating with engineers to identify and resolve problems
    \item Maintained detailed build logs, system problem reports, and traceability documentation
    \item Ensured workspace cleanliness, tool control, and compliance with aerospace manufacturing standards
    \item Trained junior team members on assembly procedures and test equipment operation
\end{itemize}}

%==================== EDUCATION ====================%
\section{Education}
\cventry{2017.03--2024.02}{\textbf{PhD in Electrical Engineering}}{Korea University}{Seoul, South Korea}{}{
\begin{itemize}[leftmargin=1cm, itemsep=0pt, parsep=0pt]
    \item Assembled and tested multi-sensor hardware through full build cycle: soldering, PCB assembly, cleanroom fabrication, and system integration
    \item Built and maintained custom test fixtures using oscilloscopes, DMMs, function generators, and power meters
    \item Performed precision cleanroom fabrication including photolithography, thin-film deposition, and chemical processing
    \item Followed iterative build-test-fix cycles across multiple hardware revisions, documenting results and process changes
\end{itemize}}

\cventry{2013.03--2017.02}{\textbf{BSc in Mechanical Engineering}}{Konkuk University}{Seoul, South Korea}{}{
\begin{itemize}[leftmargin=1cm, itemsep=0pt, parsep=0pt]
    \item Hands-on experience in CAD modeling, mechanical assembly, PCB prototyping, and precision measurement
\end{itemize}}

%==================== PUBLICATIONS ====================%
\section{Selected Publications}
\begin{sloppypar}
\begin{itemize}[leftmargin=1cm, itemsep=0pt, parsep=0pt]
    \item Jinhwee Kil, \textbf{\underline{Saeyoung Kim}}, James Jungho Pak, "Investigation of thermal management strategies for biological CubeSat payloads," \textit{Acta Astronautica}, 2024, Art. no. 228.
    \item \textbf{\underline{Kim, S.}}; Park, S.; Pak, J.J., "Multi-Modal Multi-Array Electrochemical and Optical Sensor Suite for a Biological CubeSat Payload," \textit{MDPI Sensors}, 2024, 24, Art. no. 265.
    \item \textbf{\underline{Saeyoung Kim}}, Ji-Hoon Han, Beelee Chua, James Jungho Pak, "A pH Sensing Pipette for Cross-Contamination Prevention in Industrial Fermentation," \textit{IEEE Trans. Industrial Electronics}, 2022, 69(7), 7461-7469.
\end{itemize}
\end{sloppypar}

\vfill
\centering{References available upon request.}

\end{document}
